%%%%%%%%%%%%%%%%%%%%%%%%%%%%%%%%%%%%%%%%%
% Medium Length Professional CV
% LaTeX Template
% Version 2.0 (8/5/13)
%
% This template has been downloaded from:
% http://www.LaTeXTemplates.com
%
% Original author:
% Trey Hunner (http://www.treyhunner.com/)
%
% Important note:
% This template requires the resume.cls file to be in the same directory as the
% .tex file. The resume.cls file provides the resume style used for structuring the
% document.
%
%%%%%%%%%%%%%%%%%%%%%%%%%%%%%%%%%%%%%%%%%

%----------------------------------------------------------------------------------------
%	PACKAGES AND OTHER DOCUMENT CONFIGURATIONS
%----------------------------------------------------------------------------------------

\documentclass{resume} % Use the custom resume.cls style

\usepackage[left=0.75in,top=0.6in,right=0.75in,bottom=0.6in]{geometry} % Document margins
\usepackage{hyperref} 

\name{Pavel Yakovlev}
\address{yvpavel@gmail.com}
\address{\href{https://www.yakovlev.dev}{yakovlev.dev} \\ \href{https://www.github.com/paulyakovlev}{github.com/paulyakovlev}}


\begin{document}

%----------------------------------------------------------------------------------------
%	WORK EXPERIENCE SECTION
%----------------------------------------------------------------------------------------

\begin{rSection}{Experience}
    \begin{rSubsection}{General Motors}{February 2021 - Present}{Big Data Platform Software Engineer}{Austin, TX}
    	\item Developed and published open-source contributions to Apache Hadoop projects 
    	\item Deployed and maintained large scale Hadoop components
        \item Debugged and addressed Hadoop platform issues as part of on-call rotation
        \item Develop test automation processes for build and release pipelines
    \end{rSubsection}

    \begin{rSubsection}{Lawrence Livermore National Lab}{Summer 2019}{Software Engineering Intern}{Livermore, CA}
        \item Improved automated testing and software integration framework
        \item Contributed to JUnit test suite and provided technical documentation
        \item Prototyped new test automation processess 
    \end{rSubsection}

    \begin{rSubsection}{remote.it}{May 2018 - June 2019}{Software Engineering Intern}{Palo Alto, CA}
        \item Prototyped user-end web portals for device monitoring and management
        \item Wrote a library of user-end Python and Shell scripts to manage IoT edge devices
        \item Provided technical documentation for internal and end-user software tools
    \end{rSubsection}

\end{rSection}

%----------------------------------------------------------------------------------------
%	PROJECTS SECTION
%----------------------------------------------------------------------------------------

\begin{rSection}{Projects}

    \begin{rSubsection}{Santa Cruz Stranding Map}{\href{https://github.com/lmlstrandingnetwork/lml-stranding-map}{github.com/lmlstrandingnetwork/lml-stranding-map}}{ReactJS, ExpressJS, Firebase}{}
        \item Geospatial data analysis tool for the UCSC Long Marine Lab
    \end{rSubsection}

\end{rSection}

%----------------------------------------------------------------------------------------
%	EDUCATION SECTION
%----------------------------------------------------------------------------------------

\begin{rSection}{Education}

    {\bf University of California, Santa Cruz} \hfill {\em Graduated June 2020} \\ 
    Bachelor's in Computer Science

\end{rSection}

%----------------------------------------------------------------------------------------
%	TECHNICAL SKILLS SECTION
%----------------------------------------------------------------------------------------

\begin{rSection}{Skills}
    \begin{description}
    \item[Languages:] C++, Java, Python, Javascript, HTML, CSS
    \item[Frameworks and Tools:] ReactJS, NodeJS, Flask, SQL, git, Docker, REST, Agile
    \end{description}
\end{rSection}

%----------------------------------------------------------------------------------------
%	EXAMPLE SECTION
%----------------------------------------------------------------------------------------

%\begin{rSection}{Section Name}

%Section content\ldots

%\end{rSection}

%----------------------------------------------------------------------------------------

\end{document}
